\section{The hidden quadratic twist and a rank-twelve ceiling}

The minimal-scale obstruction is a symptom of a stronger geometric failure.
For every nonzero rational $c$, the complete marked $\mu=2$ cubic-base family
has geometric generic rank at most twelve.

Put $v=t(t-1)=u^3/c$. The coefficient factors as
\[
\frac{729}{4}v^2-27(v+1)^3
=-\frac{27}{4}(v-2)^2(4v+1).
\]
Since $t$ satisfies
\[
t^2-t-\frac{u^3}{c}=0,
\]
the function field of the genus-one base is
\[
\Q(C_c)=\Q\left(u,\sqrt{1+\frac{4u^3}{c}}\right).
\]
The Mordell--Weil space therefore decomposes into the invariant curve over
$\Q(u)$ and its quadratic twist by
\[
D_c(u)=1+\frac{4u^3}{c}.
\]

\begin{theorem}[Rank-twelve ceiling for the marked cubic family]
For every $c\in\Q^*$,
\[
\rank E(\overline{\Q}(C_c))\le12.
\]
Consequently a rank-thirty specialization inside this family would require an
exceptional rank jump of at least eighteen.
\end{theorem}

\begin{proof}
The invariant surface has coefficient
\[
-\frac{27}{4}
\left(\frac{u^3}{c}-2\right)^2
\left(\frac{4u^3}{c}+1\right).
\]
It is a K3 surface with three fibres of type $IV$, three of type $II$, and an
$I_0^*$ fibre at infinity. The total root rank is
\[
3\cdot2+4=10,
\]
so Shioda--Tate and $\rho\le h^{1,1}=20$ give invariant rank at most eight.
The Euler number check is
\[
3\cdot4+3\cdot2+6=24.
\]

For a short $j=0$ model, quadratic twisting multiplies $a_6$ by $D_c^3$.
The twist coefficient is therefore
\[
-\frac{27}{4}
\left(\frac{u^3}{c}-2\right)^2
\left(\frac{4u^3}{c}+1\right)^4.
\]
This is a $\chi=3$ elliptic surface with three fibres of type $IV$ and three
of type $IV^*$. Its root rank is
\[
3\cdot2+3\cdot6=24.
\]
Since $h^{1,1}=30$, its Mordell--Weil rank is at most
\[
30-2-24=4.
\]
Its Euler number is $3\cdot4+3\cdot8=36$. The quadratic-extension rank
decomposition now gives the bound $8+4=12$.
\end{proof}

This calculation shows why the ambient Hodge capacity of a degree-three
elliptic base is not enough: repeated factors create reducible fibres whose
root lattices consume eighteen of the twenty-two hoped-for trace-zero
directions. The $\mu=2$ marked family should therefore no longer be treated
as a primary generic rank-thirty candidate.

\subsection{New setting}

The cyclic-cubic programme now moves along the CM-compatible marking conic
$q(\mu)=3s^2$ and applies the hidden-quadratic fibre-cost test before any
section search. Only parameters whose invariant and twist eigenspaces retain
combined capacity at least thirty are promoted to the expensive trisection
incidence stage.
